%! Representing a Quantitative Variable with Graphs — Unit 1, Lesson 1.3 — Warm-Up
% ONE PAGE, blank and key. Three items.
% SPOILER RULE: `histogram', `dotplot', `stemplot', `discrete', and `continuous' are what
% today's debrief attaches, so none appear here. Item 3 seeds discrete-vs-continuous in
% prior-knowledge terms; item 2 is the tallying the activity's table-fill needs.
\documentclass[10pt]{article}
\usepackage{apstats-article}
\usepackage{apstats-boxes}

\begin{document}

\pageheader{Unit 1, Lesson 1.3}{Warm-Up}

\vspace{0.15in}

\begin{enumerate}[leftmargin=1.7em, itemsep=0.19in]

  \item \textbf{From last lesson.} A theater sold 60 tickets on Friday and 240 on Saturday.
        On Friday 15 were student tickets; on Saturday 48 were.

        \vspace{5pt}
        Student share Friday: \blank{2.4cm} \qquad Saturday: \blank{2.4cm}
        \vspace{5pt}

        Which night had the larger \emph{share} of student tickets? \blank{2.8cm}

  \item Here are 12 numbers. Tally how many fall in each block.

        \vspace{4pt}
        \begin{center}
        21 \quad 34 \quad 27 \quad 45 \quad 38 \quad 22 \quad 49 \quad 31 \quad 26 \quad
        43 \quad 20 \quad 36
        \end{center}
        \vspace{4pt}

        \renewcommand{\arraystretch}{1.4}
        \begin{tabularx}{\linewidth}{@{} l Y Y Y Y @{}}
        \rowcolor{sky}
        \textbf{Block} & 20--29 & 30--39 & 40--49 & \textbf{Total} \\
        \hline
        \textbf{How many} & \blank{1.6cm} & \blank{1.6cm} & \blank{1.6cm} & \blank{1.6cm} \\
        \end{tabularx}

  \item Which of these could sensibly come out to \textbf{7.5}? \textbf{Circle} every one
        that could.
        \vspace{6pt}

        \hspace{1.0em}\textbf{(A)} the number of siblings a student has
        \hfill
        \textbf{(B)} a person's height in inches
        \hspace{1.0em}
        \vspace{5pt}

        \hspace{1.0em}\textbf{(C)} the number of cars in a parking lot
        \hfill
        \textbf{(D)} the time to run one mile, in minutes
        \hspace{1.0em}
        \vspace{7pt}

        What do the ones you circled have in common?\\[4pt]
        \writeline\\[1pt]\writeline

\end{enumerate}

\end{document}
